% GENERATED FILE. Edit canonical lessons, then run npm run generate:books.
% canonical-source-hash: fnv1a64:0d4250327c86f8ad
% canonical-lessons: MW-C125-kainchi, MW-C125-lalten, MW-C125-mombatti, MW-C125-machis, MW-C125-lakri

\chapter{Scissors, Lantern, Candle, Matches, Wood}
\label{ch:mw-scissors-lantern-candle-matches-wood}
\hlchaptermodality{125}

\begin{chapteropening}
\textbf{By the end of this chapter:} \emph{I can say scissors, a lantern, a candle, matches and wood.}

Complete the last lesson of chapter 125: I can say scissors, a lantern, a candle, matches and wood.
\end{chapteropening}

\section[\texorpdfstring{\mw{कैंची} (kaiñchī)}{कैंची (kaiñchī)}]{\mw{कैंची} (kaiñchī) — scissors}

\label{lesson:MW-C125-kainchi}



\begin{quote}
Before the new one: say the Marwadi for a needle, then the Marwadi for thread.
\end{quote}

\subsection*{You'll want to know: \mw{कैंची}}
\textbf{\mw{कैंची}} — \emph{kaiñchī} — \textquotedblleft{}scissors\textquotedblright{}.

\begin{cousinweb}
\textbf{\mw{कैंची}} (\emph{kaiñchī}) is a loan, usually traced to Persian \emph{qaičī}, \textquotedblleft{}scissors\textquotedblright{}. Persian took the word from Turkish.
\end{cousinweb}

\begin{grammarlens}[title={What you've built}]
One more word for this chapter.
\end{grammarlens}

\subsection*{Your turn}
\begin{itemize}
\raggedright
  \item \emph{Say it:} \emph{kaiñchī}
  \item \emph{Say it:} \emph{kaiñchī}, once more
  \item \emph{Say it:} say \emph{kaiñchī}
  \item \emph{Recall:} say the Marwadi for a needle, then the Marwadi for thread, then say \emph{kaiñchī} again
\end{itemize}

\begin{tcolorbox}[breakable,colback=teal!4,colframe=teal!35!black,title={Before you move on}]
What does \mw{कैंची} mean? (\textquotedblleft{}Scissors\textquotedblright{}.) Say it once more.
\end{tcolorbox}
\section[\texorpdfstring{\mw{लालटेण} (lālṭeṇ)}{लालटेण (lālṭeṇ)}]{\mw{लालटेण} (lālṭeṇ) — a lantern}

\label{lesson:MW-C125-lalten}



\begin{quote}
Before the new one: say the Marwadi for thread, then the Marwadi for scissors.
\end{quote}

\subsection*{You'll want to know: \mw{लालटेण}}
\textbf{\mw{लालटेण}} — \emph{lālṭeṇ} — \textquotedblleft{}a lantern\textquotedblright{}.

\begin{cousinweb}
\textbf{\mw{लालटेण}} (\emph{lālṭeṇ}) is a loan from English \emph{lantern}. Marwadi ends it in \emph{ṇ}.
\end{cousinweb}

\begin{grammarlens}[title={What you've built}]
One more word for this chapter.
\end{grammarlens}

\subsection*{Your turn}
\begin{itemize}
\raggedright
  \item \emph{Say it:} \emph{lālṭeṇ}
  \item \emph{Say it:} \emph{lālṭeṇ}, once more
  \item \emph{Say it:} say \emph{lālṭeṇ}
  \item \emph{Recall:} say the Marwadi for thread, then the Marwadi for scissors, then say \emph{lālṭeṇ} again
\end{itemize}

\begin{tcolorbox}[breakable,colback=teal!4,colframe=teal!35!black,title={Before you move on}]
What does \mw{लालटेण} mean? (\textquotedblleft{}A lantern\textquotedblright{}.) Say it once more.
\end{tcolorbox}
\section[\texorpdfstring{\mw{मोमबत्ती} (mombattī)}{मोमबत्ती (mombattī)}]{\mw{मोमबत्ती} (mombattī) — a candle}

\label{lesson:MW-C125-mombatti}



\begin{quote}
Before the new one: say the Marwadi for scissors, then the Marwadi for a lantern.
\end{quote}

\subsection*{You'll want to know: \mw{मोमबत्ती}}
\textbf{\mw{मोमबत्ती}} — \emph{mombattī} — \textquotedblleft{}a candle\textquotedblright{}.

\begin{cousinweb}
\textbf{\mw{मोमबत्ती}} (\emph{mombattī}) is a wax wick: \emph{mom}, \textquotedblleft{}wax\textquotedblright{}, usually traced to Persian \emph{mom}, with \emph{battī}, \textquotedblleft{}wick\textquotedblright{}, which CDIAL lists as Hindi \emph{battī} under Sanskrit \emph{varti}, \textquotedblleft{}wick\textquotedblright{} (CDIAL 11359).
\end{cousinweb}

\begin{grammarlens}[title={What you've built}]
One more word for this chapter.
\end{grammarlens}

\subsection*{Your turn}
\begin{itemize}
\raggedright
  \item \emph{Say it:} \emph{mombattī}
  \item \emph{Say it:} \emph{mombattī}, once more
  \item \emph{Say it:} say \emph{mombattī}
  \item \emph{Recall:} say the Marwadi for scissors, then the Marwadi for a lantern, then say \emph{mombattī} again
\end{itemize}

\begin{tcolorbox}[breakable,colback=teal!4,colframe=teal!35!black,title={Before you move on}]
What does \mw{मोमबत्ती} mean? (\textquotedblleft{}A candle\textquotedblright{}.) Say it once more.
\end{tcolorbox}
\section[\texorpdfstring{\mw{माचिस} (māchis)}{माचिस (māchis)}]{\mw{माचिस} (māchis) — matches}

\label{lesson:MW-C125-machis}



\begin{quote}
Before the new one: say the Marwadi for a lantern, then the Marwadi for a candle.
\end{quote}

\subsection*{You'll want to know: \mw{माचिस}}
\textbf{\mw{माचिस}} — \emph{māchis} — \textquotedblleft{}matches\textquotedblright{}.

\begin{cousinweb}
\textbf{\mw{माचिस}} (\emph{māchis}) is a loan from English \emph{matches}.
\end{cousinweb}

\begin{grammarlens}[title={What you've built}]
One more word for this chapter.
\end{grammarlens}

\subsection*{Your turn}
\begin{itemize}
\raggedright
  \item \emph{Say it:} \emph{māchis}
  \item \emph{Say it:} \emph{māchis}, once more
  \item \emph{Say it:} say \emph{māchis}
  \item \emph{Recall:} say the Marwadi for a lantern, then the Marwadi for a candle, then say \emph{māchis} again
\end{itemize}

\begin{tcolorbox}[breakable,colback=teal!4,colframe=teal!35!black,title={Before you move on}]
What does \mw{माचिस} mean? (\textquotedblleft{}Matches\textquotedblright{}.) Say it once more.
\end{tcolorbox}
\section[\texorpdfstring{\mw{लकड़ी} (lakṛī)}{लकड़ी (lakṛī)}]{\mw{लकड़ी} (lakṛī) — wood}

\label{lesson:MW-C125-lakri}



\begin{quote}
Before the new one: say the Marwadi for a candle, then the Marwadi for matches.
\end{quote}

\subsection*{You'll want to know: \mw{लकड़ी}}
\textbf{\mw{लकड़ी}} — \emph{lakṛī} — \textquotedblleft{}wood\textquotedblright{}.

\begin{cousinweb}
\textbf{\mw{लकड़ी}} (\emph{lakṛī}) is a tadbhava, a word inherited from Sanskrit and reshaped by sound change. It comes from \emph{lakuṭa}, \textquotedblleft{}club\textquotedblright{}. CDIAL gives no Marwari form in this entry (CDIAL 10875); the match is through Hindi \emph{lakṛī}.
\end{cousinweb}

\begin{grammarlens}[title={What you've built}]
One more word for this chapter.
\end{grammarlens}

\subsection*{Your turn}
\begin{itemize}
\raggedright
  \item \emph{Say it:} \emph{lakṛī}
  \item \emph{Say it:} \emph{lakṛī}, once more
  \item \emph{Say it:} say \emph{lakṛī}
  \item \emph{Recall:} say the Marwadi for a candle, then the Marwadi for matches, then say \emph{lakṛī} again
\end{itemize}

\begin{tcolorbox}[breakable,colback=teal!4,colframe=teal!35!black,title={Before you move on}]
What does \mw{लकड़ी} mean? (\textquotedblleft{}Wood\textquotedblright{}.) Say it once more.
\end{tcolorbox}
